\section*{Conventions and scope}
We write $F_2=\langle x,y\rangle$, $z=(xy)^{-1}$, and
$[a,b]=aba^{-1}b^{-1}$. A hat on a discrete group means its full profinite
completion; $\SL_2(\Zhat)=\prod_p\SL_2(\Z_p)$ means its congruence completion.
These two operations are never interchanged. Homomorphisms of profinite groups
are continuous. Arithmetic fundamental groups use a geometric or tangential
base point; changing a path gives the indicated inner-conjugacy ambiguity.
We write $G_\Q=\Gal(\overline\Q/\Q)$ and $\chi$ for the cyclotomic character.
The symbols in the manuscript are fixed in the dictionary in Part A; primes, exclamation marks, and tildes are mathematical data.
The primary manuscript is cited as \citet{Grothendieck1981root}.

\noindent\textbf{Reading conventions.}
A reference ``140-3:209, f.\ 403 verso'' means the archive cote, scan, and handwritten
folio; where folio attribution is not unambiguous, we retain the scan as the
stable locator and explicitly say so. Corrected transcriptions are used as
working text, with the selected original scans inspected whenever a conclusion
depends on a disputed formula. Mathematical consistency provides editorial
evidence, not a final palaeographic reading. Withdrawn assertions are marked as
withdrawn. The notes, LM3 return, and reviews are leads rather than authority.

\noindent\textbf{Status conventions.}
\textsf{Proved} includes a proof here; \textsf{Known} includes an exact source
and the reduction required in this setting; \textsf{Refuted} includes an explicit
counterexample and a corrected assertion; \textsf{Open} records the remaining
precise assertion and the attack. Conditional use of an unrefereed preprint is
named as such. A bibliography entry marked verified means that we opened the
source; it does not mean that we independently verified its entire proof.
The report distinguishes partly settled tasks from unfinished ones whenever
the requested audit could not reach its stated standard.
